\ifdefined\gkver\else\def\gkver{student}\fi
\documentclass[\gkver]{gaokaozhenti}
\begin{document}

\chapter{2003年广东综合}
%% paperType: 综合能力测试物理部分

\begin{enumerate}
\item
%% number: 31
%% paperName: 2003年广东综合第 31 题 6 分
%% typeId: 1
%% type: 单选题
%% chapter: 功和能
%% point: 动能定理
%% method: 无
%% score: 6
%% degree: 750
%% duplicateId: 0
%% body:
在离地面高为 $h$ 处竖直上抛一质量为 $m$ 的物块，抛出时的速度为 $v_{0}$，当它落到地面时速度为 $v$，用 $g$ 表示重力加速度，则在此过程中物块克服空气阻力所做的功等于\xzanswer{C}

\fourchoices[answer=C]
{$mgh - \frac{1}{2}mv^{2} - \frac{1}{2}mv_{0}^{2}$}
{$- \frac{1}{2}mv^{2} - \frac{1}{2}mv_{0}^{2} - mgh$}
{$mgh + \frac{1}{2}mv_{0}^{2} - \frac{1}{2}mv^{2}$}
{$mgh + \frac{1}{2}mv^{2} - \frac{1}{2}mv_{0}^{2}$}

%% answer: C
%% memo:
\memoanswer{
本题原卷及材料中未提供详解，待补充。
}

\item
%% number: 32
%% paperName: 2003年广东综合第 32 题 6 分
%% typeId: 1
%% type: 单选题
%% chapter: 曲线运动
%% point: 人造地球卫星
%% method: 无
%% score: 6
%% degree: 850
%% duplicateId: 0
%% body:
若航天飞机在一段时间内保持绕地心做匀速圆周运动，则\xzanswer{C}

\fourchoices[answer=C]
{它的速度的大小不变，动量也不变}
{它不断地克服地球对它的万有引力做功}
{它的动能不变，引力势能也不变}
{它的速度的大小不变，加速度等于零}

%% answer: C
%% memo:
\memoanswer{
本题原卷及材料中未提供详解，待补充。
}

\item
%% number: 33
%% paperName: 2003年广东综合第 33 题 6 分
%% typeId: 1
%% type: 单选题
%% chapter: 磁场
%% point: 带电粒子在复合场中的运动
%% method: 无
%% score: 6
%% degree: 750
%% duplicateId: 0
%% body:
如右图所示，a、b 是位于真空中的平行金属板，a 板带正电，b 板带负电，两板间的电场为匀强电场，场强为 $E$。同时在两板之间的空间中加匀强磁场，磁场方向垂直于纸面向里，磁感应强度为 $B$。一束电子以大小为 $v_{0}$ 的速度从左边 S 处沿图中虚线方向入射，虚线平行于两板，要想使电子在两板间能沿虚线运动，则 $v_{0}$、$E$、$B$ 之间的关系应该是\xzanswer{A}

\onepicture[w=5.3cm]{figs/33.png}

\fourchoices[answer=A]
{$v_{0} = \frac{E}{B}$}
{$v_{0} = \frac{B}{E}$}
{$v_{0} = \sqrt{\frac{E}{B}}$}
{$v_{0} = \sqrt{\frac{B}{E}}$}

%% answer: A
%% memo:
\memoanswer{
本题原卷及材料中未提供详解，待补充。
}

\item
%% number: 34
%% paperName: 2003年广东综合第 34 题 6 分
%% typeId: 1
%% type: 单选题
%% chapter: 电磁感应
%% point: 法拉第电磁感应定律
%% method: 无
%% score: 6
%% degree: 650
%% duplicateId: 0
%% body:
电学中的库仑定律、欧姆定律、法拉第电磁感应定律（有关感应电动势大小的规律）、安培定律（磁场对电流作用的规律）都是一些重要的规律，右图为远距离输电系统的示意图（为了简单，设用户的电器是电动机），下列选项中正确的是\xzanswer{D}

\onepicture[w=9.5cm]{figs/34.png}

\fourchoices[answer=D]
{发电机能发电的主要原理是库仑定律，变压器能变压的主要原理是欧姆定律，电动机通电后能转动起来的主要原理是法拉第电磁感应定律}
{发电机能发电的主要原理是安培定律，变压器能变压的主要原理是欧姆定律，电动机通电后能转动起来的主要原理是库仑定律}
{发电机能发电的主要原理是欧姆定律，变压器能变压的主要原理是库仑定律，电动机通电后能转动起来的主要原理是法拉第电磁感应定律}
{发电机能发电的主要原理是法拉第电磁感应定律，变压器能变压的主要原理是法拉第电磁感应定律，电动机通电后能转动起来的主要原理是安培定律}

%% answer: D
%% memo:
\memoanswer{
本题原卷及材料中未提供详解，待补充。
}

\item
%% number: 35
%% paperName: 2003年广东综合第 35 题 6 分
%% typeId: 1
%% type: 单选题
%% chapter: 电路
%% point: 串并联电路
%% method: 无
%% score: 6
%% degree: 750
%% duplicateId: 0
%% body:
在右图所示电路中 $E$ 为电源，其电动势为 $9.0\UV$，内阻可忽略不计；AB 为滑动变阻器。其电阻 $R = 30\UO$；L 为一小灯泡，其额定电压 $U = 6.0\UV$，额定电阻 $P = 1.8\UW$；K 为电键。开始时滑动变阻器的触头位于 B 端，现在接通电键 K，然后将触头缓慢地向 A 端滑动，当到达某一位置 C 处时，小灯泡刚好正常发光，则 CB 之间的电阻应为\xzanswer{B}

\onepicture[w=5.3cm]{figs/35.png}

\fourchoices[answer=B]
{$10\UO$}
{$20\UO$}
{$15\UO$}
{$5\UO$}

%% answer: B
%% memo:
\memoanswer{
本题原卷及材料中未提供详解，待补充。
}

\item
%% number: 36
%% paperName: 2003年广东综合第 36 题 6 分
%% typeId: 1
%% type: 单选题
%% chapter: 运动和力
%% point: 牛顿第二定律
%% method: 整体与隔离
%% score: 6
%% degree: 450
%% duplicateId: 0
%% body:
如右图所示，一质量为 $M$ 的楔形木块放在水平桌面上，它的顶角为 $90^{\circ}$，两底角为 $\alpha$ 和 $\beta$；a、b 为两个位于斜面上质量均为 $m$ 的小木块。已知所有接触面都是光滑的。现发现 a、b 沿斜面下滑，而楔形木块静止不动，这时楔形木块对水平桌面的压力等于\xzanswer{A}

\onepicture[w=6.3cm]{\includesvg{figs/36}}

\fourchoices[answer=A]
{$Mg + mg$}
{$Mg + 2mg$}
{$Mg + mg(\sin\alpha + \sin\beta)$}
{$Mg + mg(\cos\alpha + \cos\beta)$}

%% answer: A
%% memo:
\memoanswer{
本题原卷及材料中未提供详解，待补充。
}

\item
%% number: 37
%% paperName: 2003年广东综合第 37 题 14 分
%% typeId: 1
%% type: 单选题
%% chapter: 光学
%% point: 光的电磁说
%% method: 无
%% score: 14
%% degree: 650
%% duplicateId: 0
%% body:
二氧化碳能强烈吸收红外长波辐射，这种长波辐射的波长范围约是 $1.4\times10^{-3}\sim6\times10^{-3}$，相应的频率范围是\tkanswer[2.0cm]{}，相应的光子能量范围是\tkanswer[2.0cm]{}。“温室效应”使大气全年的平均温度升高，空气温度升高，从微观上看就是空气中分子的\tkanswer[2.2cm]{}。（已知普朗克恒量 $h=6.6\times10^{-34}\UJcdots$，真空中的光速 $c=3.0\times10^{8}\Umdsn$。结果取两位数字。）
%% answer: H
\jdanswer{
$2.1\times10^{12}\sim1.9\times10^{13}\UHz$；$1.4\times10^{-20}\sim1.3\times10^{-20}\UJ$；无规则运动（或热运动）更剧烈、平均动能更大。
}
%% memo:
\memoanswer{
本题原卷及材料中未提供详解，待补充。
}

\end{enumerate}

\end{document}
